\chapter{Arquitetura e ambiente de testes}
\label{cap2}

Este capítulo descreve o ambiente laboratorial utilizado para o desenvolvimento e validação do projeto, incluindo a infraestrutura virtualizada, a organização cliente-servidor, a stack técnica adotada e os principais critérios de teste. O objetivo deste ambiente foi garantir condições controladas para instalar, configurar e verificar de forma automatizada os serviços de Administração de Sistemas previstos no enunciado.

O laboratório foi implementado em Oracle VM VirtualBox, sobre um computador anfitrião com sistema operativo Windows, que também foi utilizado como plataforma complementar de teste em algumas etapas. A opção por virtualização permitiu repetir configurações, isolar o contexto de execução e reduzir riscos sobre o sistema anfitrião, mantendo simultaneamente flexibilidade para reconfiguração durante o desenvolvimento.

Foram utilizadas duas máquinas virtuais principais, ambas com AlmaLinux 8. A máquina \textit{AS\_Server} foi configurada como servidor central do projeto, com disco principal de 20 GB e três discos adicionais de 5 GB destinados à implementação de RAID 5. A máquina \textit{AS\_Client}, também com AlmaLinux 8, e disco de 20 GB, foi utilizada para testar os serviços disponibilizados pelo servidor. Em ambos os casos foi selecionado no VirtualBox o tipo Linux e o subtipo Red Hat (64-bit)

A conetividade entre as máquinas foi definida em modo \textit{Bridged Adapter}, esta escolha permitiu que as máquinas virtuais se comportassem como nós autónomos na rede local, com comunicação direta entre cliente e servidor e acesso à Internet para instalação de pacotes e dependências. Importa salientar que o endereçamento IP não foi fixo e variaram consoante a rede onde o laboratório era executado.

Após a instalação base do sistema operativo nas duas VMs, foram aplicadas configurações iniciais de preparação do ambiente: atualização do sistema, instalação de pacotes essenciais, ajuste de hostname, preparação de acesso remoto por SSH e criação da diretoria de trabalho para os scripts (\texttt{/root/projeto}). Para acelerar a fase inicial de validação laboratorial e reduzir interferências externas, foram também adotadas medidas de simplificação temporária, nomeadamente desativação de SELinux e da \textit{firewalld}.

Do ponto de vista de implementação, o projeto foi desenvolvido  em Bash, por oferecer integração direta com ferramentas nativas de administração Linux e boa adequação à automatização de tarefas sequenciais. A stack técnica incluiu os serviços e utilitários necessários para cobrir os requisitos do enunciado: BIND/named para DNS, Apache httpd para serviço Web e VirtualHosts, Samba para partilhas com clientes Windows, NFS para partilhas com clientes Linux/Unix, \texttt{tar} e \texttt{rsync} para backups, \texttt{mdadm} para RAID 5, Fail2Ban para mitigação de ataques de força bruta a SSH. 

Em síntese, a arquitetura adotada privilegiou simplicidade, repetibilidade e alinhamento com os requisitos do projeto. Embora o ambiente tenha sido desenhado para fins laboratoriais e académicos, com algumas simplificações controladas de segurança em fases iniciais, a sua estrutura permitiu suportar com eficácia o desenvolvimento dos scripts, a demonstração dos serviços e a produção de evidências técnicas para o relatório final.


